% !TeX root = ../main.tex
\chapter{تحلیل وابستگی و کشف زیرگروه}
\source{اسلاید تحلیل وابستگی، PDF ۳ تا ۴۶؛ کتاب فصل ۴، صفحات چاپی ۲۱۳ تا ۲۷۴ و بخش‌های ۷٫۱ و ۷٫۲. FP-growth و مجموعه‌های بسته بسط کتاب‌اند؛ WRAcc و کشف زیرگروه بر اسلایدها تکیه دارند.}
\section{سؤال درست و دادهٔ تراکنشی}

تحلیل ارتباط می‌پرسد «کدام رخدادها در یک واحد مشاهده کنار هم ظاهر می‌شوند؟» واحد می‌تواند یک سبد خرید، یک نشست وب یا مجموعهٔ علائم یک بیمار باشد. ابتدا باید تعریف تراکنش روشن شود؛ تجمیع همهٔ خریدهای یک مشتری با جدا نگه داشتن هر فاکتور، پشتیبانی و اطمینان متفاوتی می‌دهد (تمرین۲ کتاب چاپی۲۹۵/PDF ۲۹۶). هر تراکنش مجموعه است؛ تکرار چند واحد شیر در مدل دودویی یک حضور شمرده می‌شود. صفرِ نخریدن هزاران کالای نامرتبط معمولاً از یکِ خریدن اهمیت کمتری دارد، پس داده دودویی نامتقارن است.

اگر $I$ مجموعهٔ $d$ قلم و $T$ شامل $N$ تراکنش باشد، برای $X\subseteq I$ داریم:

\[\sigma(X)=\sum_{i=1}^{N}\mathbf1[X\subseteq t_i],\qquad s(X)=\sigma(X)/N.\]

دو سوی قاعدهٔ $X\to Y$ باید ناتهی و مجزا باشند: $X\cap Y=\varnothing$. تعریف پشتیبانی و اطمینان چنین است:
\[s(X\to Y)=s(X\cup Y),\qquad c(X\to Y)=\frac{s(X\cup Y)}{s(X)}.\]
اطمینان، احتمال شرطی تجربیِ تالی در تراکنش‌های دارای مقدم است؛ با پشتیبانی مقدم اشتباه نشود. اگر پشتیبانی مقدم صفر باشد، اطمینان تعریف ندارد. پیکان، علت یا ترتیب زمانی را ثابت نمی‌کند. همبستگی دو ستون دودویی همان ضریب فی است و اگر یکی ثابت باشد تعریف‌نشده می‌شود.

\section{مثال عددی مرجع، یکسان با اسلاید و کتاب}

برای رفع خطای املایی \lr{Break} در برخی اسلایدها، همه‌جا \lr{Bread} منظور است. نمادهای کوتاه این مثال عبارت‌اند از: $B$ نان، $M$ شیر، $D$ پوشک، $R$ نوشیدنی، $E$ تخم‌مرغ و $C$ نوشابه. در این نمایش آموزشی، $R$ معادل قلم نوشیدنی منبع است.

{\def\LTcaptype{table} % do not increment counter
\begin{longtable}[]{@{}rl@{}}
\toprule\noalign{}
تراکنش & اقلام \\
\midrule\noalign{}
\endhead
\bottomrule\noalign{}
\endlastfoot
۱ & $\{B,M\}$ \\\
۲ & $\{B,D,R,E\}$ \\\
۳ & $\{M,D,R,C\}$ \\\
۴ & $\{B,M,D,R\}$ \\\
۵ & $\{B,M,D,C\}$ \\\
\end{longtable}
}

شمارش تک‌قلم‌ها به‌ترتیب چنین است:
\[\sigma(B)=4,\quad\sigma(M)=4,\quad\sigma(D)=4,\quad\sigma(R)=3,\quad\sigma(C)=2,\quad\sigma(E)=1.\]

برای قاعدهٔ $\{M,D\}\to\{R\}$، شمارش مقدم ۳، شمارش تالی ۳ و شمارش اجتماع ۲ است. بنابراین:
\begin{align*}
s(\{M,D\}\to\{R\})&=\frac25=0.4,\\
c(\{M,D\}\to\{R\})&=\frac23\approx0.6667,\\
\operatorname{Lift}(\{M,D\}\to\{R\})&=\frac{2/3}{3/5}=\frac{10}{9}\approx1.1111.
\end{align*}
قانون معکوس $\{R\}\to\{M,D\}$ نیز اطمینان $2/3$ دارد، اما این برابری اتفاقی است و اطمینان به‌طور کلی متقارن نیست.

برای قاعدهٔ $D\to R$ داریم:
\begin{align*}
s(D\to R)&=\frac35=0.6,\\
c(D\to R)&=\frac34=0.75,\\
\operatorname{Lift}(D\to R)&=\frac{0.75}{0.6}=1.25.
\end{align*}
در جهت $R\to D$ اطمینان یک است؛ بالابری همان $1.25$ باقی می‌ماند. در نتیجه اطمینان جهت‌دار و بالابری متقارن است.

\section{چرا جست‌وجوی خام جواب نمی‌دهد؟}

برای $d$ قلم، $2^d-1$ زیرمجموعهٔ ناتهی داریم. از یک مجموعهٔ $k$عضوی، $2^k-2$ قاعده با دو سوی ناتهی ساخته می‌شود. تعداد تمام قواعد ممکن برابر است با
\[3^d-2^{d+1}+1.\]
هر قلم در مقدم، تالی یا هیچ‌کدام است؛ سپس حالت‌های دارای یک سوی تهی کنار گذاشته می‌شوند. برای شش قلم، ۶۰۲ قاعده ممکن است. تطبیق خام $M$ نامزد با $N$ تراکنش و پهنای حداکثر $w$ حدود $O(NMw)$ هزینه دارد. رشد نمایی، جست‌وجوی خام را حتی برای دادهٔ متوسط دشوار می‌کند.

\section{Apriori ، منطق و اجرای واقعی}

اگر $X\subseteq Y$ باشد، هر تراکنش حاوی $Y$ حتماً $X$ را هم دارد، پس $s(Y)\le s(X)$. بنابراین اگر $X$ کم‌تکرار است، تمام ابرمجموعه‌هایش را می‌توان بدون شمارش حذف کرد. عکس گزاره برقرار نیست: پرتکرار بودن همهٔ زیرمجموعه‌ها تضمین پرتکرار بودن اجتماع نیست.

الگوریتم: تک‌اقلام را بشمار؛ مجموعهٔ پرتکرار $F_1$ را بساز؛ از $F_k$ نامزد $C_{k+1}$ با اتصال پیشوند مشترک بساز؛ هر نامزدی را که یک زیرمجموعهٔ $k$ عضوی کم‌تکرار دارد حذف کن؛ با اسکن داده پشتیبانی بازمانده‌ها را بشمار؛ تکرار تا خالی شدن $F$. ترتیب واژه‌نامه‌ای از نامزد تکراری جلوگیری می‌کند. حمایت ضدیکنوا شرط هرس است، نه «اعتماد به مشابه بودن نام‌ها».

با حداقل شمارش ۳، یعنی پشتیبانی ۶۰ درصد، مراحل برای دادهٔ مرجع چنین است:
\begin{enumerate}
\item تک‌اقلام پرتکرار $F_1=\{B,M,D,R\}$ هستند؛ $C,E$ حذف می‌شوند.
\item شمارش شش زوج: $BM=3,BD=3,BR=2,MD=3,MR=2,DR=3$.
\item بنابراین $F_2=\{BM,BD,MD,DR\}$ است.
\item تنها نامزد سه‌تایی $BMD$ است، چون همهٔ زوج‌هایش پرتکرارند. شمارش واقعی آن ۲ است و حذف می‌شود.
\item $F_3$ تهی و فرایند تمام می‌شود. در این اجرا $6+6+1=13$ نامزد شمرده شده، درحالی‌که تعداد مجموعه‌های ناتهی تا اندازهٔ سه ۴۱ و تمام اندازه‌ها ۶۳ است.
\end{enumerate}

حداقل شمارش از پشتیبانی نسبی باید $\lceil N\times\mathrm{minsup}\rceil$ باشد؛ مثلاً $N=5$ و $\mathrm{minsup}=0.5$، حداقل شمارش۳ است. مرز «بزرگ‌تر یا مساوی» را یکسان رعایت کنیم. کاهش minsup ، افزایش تعداد اقلام، تعداد تراکنش‌ها و پهنای سبد می‌توانند هزینه را بالا ببرند. در دادهٔ متراکم الگوهای بلند بسیار زیادند و هیچ الگوریتمی هزینهٔ خروجی نمایی را جادویی حذف نمی‌کند.

\section{تولید قواعد و هرس درست اطمینان}

از هر مجموعهٔ پرتکرار $L$، افرازهای $X\to L\setminus X$ تولید می‌شوند. شمارش زیرمجموعه‌ها از مرحلهٔ قبل موجود است. اطمینان در تغییر آزادِ مقدم، ضدیکنوا نیست؛ اما با ثابت ماندن اجتماع $L$، بزرگ‌تر کردن تالی، مقدم را کوچک‌تر و پشتیبانی مقدم را بزرگ‌تر یا مساوی می‌کند. صورت کسر ثابت است، پس اطمینان افزایش نمی‌یابد. در نتیجه، اگر یک تالی آستانهٔ اطمینان را نگذرد، تالی‌های بزرگ‌تر شامل آن قابل هرس‌اند. این حکم را نمی‌توان به قواعدی با اجتماع متفاوت تعمیم داد.

نمونهٔ $L=\{M,D,R\}$ را در آستانهٔ پشتیبانی ۴۰ درصد در نظر بگیرید. اطمینان شش قاعده چنین است:
\begin{align*}
c(\{M,D\}\to R)&=\frac23,& c(\{M,R\}\to D)&=1,\\
c(\{D,R\}\to M)&=\frac23,& c(M\to\{D,R\})&=\frac12,\\
c(D\to\{M,R\})&=\frac12,& c(R\to\{M,D\})&=\frac23.
\end{align*}
اگر $\mathrm{minconf}=0.7$ باشد، فقط $\{M,R\}\to D$ موفق است؛ نباید $L$ را با آستانهٔ ۶۰ درصد مثال قبل پرتکرار معرفی کرد.

\section{جذابیت، نرخ پایه و Lift}

اطمینان بالا به‌تنهایی همراهی مثبت را نشان نمی‌دهد. معیار بالابری یا Lift ، همراهی را نسبت به نرخ پایه می‌سنجد:
\[\operatorname{Lift}(X\to Y)=\frac{c(X\to Y)}{s(Y)}=\frac{s(X\cup Y)}{s(X)s(Y)}.\]
مقدار بیشتر از یک همراهی مثبت تجربی، مقدار یک استقلال در جدول تجربی و مقدار کمتر از یک همراهی منفی را بیان می‌کند. این معیار آزمون معنی‌داری یا سند علّیت نیست؛ چند مشاهدهٔ نادر می‌توانند مقدار بسیار بالا و ناپایداری تولید کنند.

\textbf{مثال اسلایدهای ۳۲ تا ۳۵.} جدول چای و قهوه، با ترتیب خانه‌های $[a,b;c,d]$، شامل $[3,1;15,1]$ و جمعیت $N=20$ است. برای قاعدهٔ چای به قهوه داریم:
\begin{align*}
P(\mathrm{coffee})&=\frac{18}{20}=0.9,\\
c(\mathrm{tea}\to\mathrm{coffee})&=\frac34=0.75,\\
s(\mathrm{tea}\to\mathrm{coffee})&=\frac3{20}=0.15,\\
\operatorname{Lift}&=\frac{0.75}{0.9}=\frac56\approx0.8333.
\end{align*}
دانستن چای‌نوش بودن احتمال قهوه را از ۹۰ درصد به ۷۵ درصد کاهش می‌دهد.

\textbf{مثال مستقل کتاب، صفحات چاپی ۲۵۹ تا ۲۶۱.} جدول $[150,50;650,150]$، با $N=1000$، نرخ پایهٔ قهوه ۸۰ درصد و اطمینان ۷۵ درصد دارد؛ بالابری $0.9375$ است. این را با اعداد اسلاید مخلوط نکنید. مثال چای و عسل کتاب، جدول $[100,100;20,780]$، نرخ پایهٔ عسل ۱۲ درصد، اطمینان ۵۰ درصد و بالابری تقریباً $4.1667$ دارد: اطمینان کمتر ولی ارتباط مثبت قوی‌تر.

افزونگی: اگر افزودن شرطی به مقدم هیچ پوشش/اطمینان جدیدی ندهد، قاعدهٔ طولانی‌تر ممکن است اطلاعات افزوده نداشته باشد. معیار عددی با فایدهٔ عملی، تازگی و دانش دامنه تکمیل می‌شود؛ بالاترین Lift همیشه بهترین تصمیم نیست.

\section{بسط کنترل‌شدهٔ کتاب: نمایش فشرده و FP-growth}

مجموعهٔ پرتکرار بیشینه هیچ ابرمجموعهٔ پرتکراری ندارد؛ مجموعهٔ بسته هیچ ابرمجموعهٔ دارای همان پشتیبانی ندارد. همهٔ مجموعه‌های پرتکرار بیشینه بسته‌اند، ولی عکس درست نیست. بیشینه‌ها وجود الگوهای پرتکرار کوچک‌تر را بیان می‌کنند اما شمارش پشتیبانی آن‌ها را حفظ نمی‌کنند؛ مجموعه‌های بسته همراه شمارش می‌توانند شمارش تمام مجموعه‌های پرتکرار را بازسازی کنند.

در نمایش عمودی برای هر قلم فهرست شناسهٔ تراکنش‌ها را نگه می‌داریم و اشتراک فهرست‌ها شمارش اجتماع را می‌دهد. این به Jaccard شباهت ظاهری دارد اما هدف الگوریتم شمارش پشتیبانی است.

FP-growth به جای تولید نامزدهای همهٔ سطوح، داده را در FP-tree فشرده می‌کند: گذر اول شمارش و حذف تک‌اقلام کم‌تکرار؛ مرتب‌سازی مشترک اقلام باقی‌مانده (معمولاً نزولی پشتیبانی با شکست تساوی ثابت)؛ گذر دوم درج هر سبد در مسیر درخت و افزایش شمارندهٔ پیشوند مشترک. جدول سرآیند و پیوند گره‌های هم‌نام دسترسی به مسیرها را آسان می‌کند. سپس برای هر پسوند، مسیرهای پیشوندی با \textbf{وزن شمارندهٔ گره همان پسوند} استخراج می‌شوند؛ این پایگاه الگوی شرطی است. درخت شرطی با حذف قلم پسوند و اقلام کم‌تکرار ساخته و فرایند بازگشتی تکرار می‌شود. شمارندهٔ والد را نباید به اشتباه وزن همهٔ مسیرهای شرطی گرفت. مسیر تهی برای ساخت الگوی بلندتر چیزی اضافه نمی‌کند. شمارندهٔ یک گره لزوماً پشتیبانی کلی قلمش نیست؛ ممکن است چند گره با همان برچسب وجود داشته باشد. فشرده‌سازی و کارایی به اشتراک پیشوندها وابسته است؛ مرتب‌سازی نزولی الزاماً کمینه‌ترین درخت ممکن را تضمین نمی‌کند. این بخش بسط تکمیلی کتاب است.


\subsection*{پایگاه شرطی را واقعاً بسازیم}
در همان پنج سبد و حداقل شمارش ۳، اقلام $C,E$ حذف می‌شوند. برای شکست تساوی، ترتیب مشترک $B,D,M,R$ را ثابت بگیرید؛ فراوانی سه قلم اول برابر ۴ و قلم آخر ۳ است. سبدهای مرتب، به‌ترتیب $BM$، $BDR$، $DMR$، $BDMR$ و $BDM$ هستند. مسیرهای منتهی به $R$، پس از حذف خود این قلم، پایگاه شرطی زیر را می‌دهند:
\[ (BD,1),\qquad(BDM,1),\qquad(DM,1).\]
وزن‌ها از شمارندهٔ گرهٔ $R$ هر مسیر آمده‌اند. در این پایگاه، شمارش‌های شرطی $B=2,D=3,M=2$ هستند؛ با حداقل شمارش ۳ فقط $D$ می‌ماند. بنابراین الگوی همراه $R$ که از این شاخه به دست می‌آید $DR$ با شمارش ۳ است. حذف $R$ از مسیر، به معنی حذف وزن تراکنش‌های منتهی به آن نیست. این اجرای کوتاه، تفاوت «شمارندهٔ یک گره» و «پشتیبانی کل یک قلم» را نیز روشن می‌کند.

\section{شباهت و معیارهای تکمیلی}

برای جدول دو متغیر دودویی، $a=f_{11},b=f_{10},c=f_{01},d=f_{00}$ و $N=a+b+c+d$ است. چند معیار پرکاربرد عبارت‌اند از
\begin{align*}
J&=\frac a{a+b+c},&\operatorname{cos}&=\frac a{\sqrt{(a+b)(a+c)}},\\
\phi&=\frac{ad-bc}{\sqrt{(a+b)(a+c)(b+d)(c+d)}},&\operatorname{Lift}&=\frac{Na}{(a+b)(a+c)}.
\end{align*}
اختلاف همراهی از استقلال نیز $\operatorname{leverage}=a/N-[(a+b)/N][(a+c)/N]$ است. در هر معیار، حالت مخرج صفر باید جداگانه مدیریت شود.

\lr{Jaccard} و \lr{cosine} با افزودن تراکنش‌هایی که هیچ‌کدام از دو قلم را ندارند تغییر نمی‌کنند، ولی بالابری تغییر می‌کند. در مثال کتاب، صفحات چاپی ۲۶۹ تا ۲۷۰، جدول $[700,100;100,100]$ را به $[700,100;100,1100]$ تبدیل می‌کنیم. نتایج عبارت‌اند از:
\begin{align*}
\operatorname{cos}&=0.875,& J&=\frac79,\\
\operatorname{Lift}_{\mathrm{before}}&=1.09375,&
\operatorname{Lift}_{\mathrm{after}}&=2.1875.
\end{align*}
\textbf{خود متن کتاب در این مثال، بالابری‌های $0.972$ و $1.944$ را چاپ کرده که با جدول ناسازگارند؛ از محاسبهٔ مستقیم استفاده شود.} مخرج صفر برای هر معیار باید مدیریت شود. شباهت بدون توجه به صفرهای مشترک برای سبد خرید معمولاً معنادارتر است؛ برای ویژگی دودویی متقارن، صفر نیز اطلاعات دارد.

پارادوکس سیمپسون (بسط، چاپی۲۷۲--۲۷۴/PDF ۲۷۳--۲۷۵): در دادهٔ تجمیعی خریداران تلویزیون۵۵٪ و غیرخریداران۴۵٪ ابزار ورزشی می‌خرند؛ اما بین دانشجویان $1/10<4/34$ و بین شاغلان $98/170<50/86$ است. ترکیب متفاوت گروه‌ها علامت ارتباط را عوض می‌کند. تفسیر رابطه باید متغیرهای مخدوش‌کننده و تعریف جامعه را در نظر بگیرد.

\section{ویژگی‌های اسمی، پیوسته و کشف زیرگروه}

ویژگی اسمی به اقلام «نام ویژگی=مقدار» تبدیل می‌شود. ویژگی تک‌مقداری مانند مرورگر نمی‌تواند هم‌زمان دو مقدار داشته باشد؛ نامزدهای شامل دو مقدار ناسازگار همان ویژگی را از ابتدا حذف کنیم. مقادیر نادر را فقط با توجیه معنایی یا گروه «سایر» تجمیع کنیم. قلم بسیار رایج ممکن است قواعد بدیهی بسازد؛ حذف کورکورانهٔ آن می‌تواند رابطهٔ معنادار را هم از بین ببرد، پس هدف و نرخ پایه مهم‌اند.

ویژگی پیوسته را می‌توان با بازه‌های هم‌عرض یا هم‌فراوانی گسسته کرد. بازهٔ خیلی باریک حمایت کافی ندارد؛ بازهٔ خیلی پهن ممکن است گروه‌های متفاوت را مخلوط و اطمینان را تضعیف کند. مرز و شمول ابتدا/انتها روشن باشد، مثلاً $[5,10)$. مثال کتاب با سن کاربران و چت: گروه۱۶--۲۴ شامل۲۲ پاسخ مثبت از۲۷ نفر در کل۲۵۰ مشاهده است؛ $s=22/250=0.088$ و $c=22/27\approx0.8148$. شکستن به۱۶--۲۰ و۲۰--۲۴ هرکدام حمایت $11/250=0.044$ دارد؛ ادغام در بازهٔ خیلی پهن۱۲--۳۶ اطمینان $75/130\approx0.5769$ می‌دهد. اثر مرزبندی را می‌توان در آزمایشگاه بررسی کرد.

کشف زیرگروه میان توصیف و پیش‌بینی است: هدف، یافتن توصیف‌های قابل فهم از گروه‌هایی با رفتار متفاوتِ هدف $T$ است، نه فقط کمینه‌کردن خطای یک طبقه‌بند جامع. زیرگروه‌ها می‌توانند هم‌پوشان باشند و لزومی ندارد همهٔ جمعیت را بپوشانند. در اسلاید۴۱ توصیف کوتاه «خرطوم دارد» و توصیف طولانی‌تر فیل مطرح می‌شود؛ برای پیش‌بینی سادگی و تعمیم مهم است، برای شناخت داده توصیف‌های متعدد می‌توانند مفید باشند. الگوریتم‌های مطرح‌شده: استخراج قواعد از درخت هرس‌نشده و رتبه‌بندی (EXPLORA/MIDOS)، روش‌های مبتنی بر Apriori-SD ، روش‌های تکاملی. این روش‌ها مسیرهای متفاوت جست‌وجوی زیرگروه‌اند؛ امتیاز و فضای توصیف، نتیجه را تعیین می‌کند.

اگر $S$ زیرگروه و $T$ رویداد هدف باشد، معیار صحت نسبی وزن‌دار چنین است:
\[\operatorname{WRAcc}(S,T)=P(S\cap T)-P(S)P(T)=P(S)\{P(T\mid S)-P(T)\}.\]
امتیاز هم پوشش و هم افزایش نرخ هدف نسبت به پایه را می‌سنجد. گروه بسیار کوچک و خالص ممکن است امتیاز کمتری از گروه بزرگ‌تر با خلوص متوسط بگیرد. استقلال امتیاز صفر و نرخ کمتر از پایه امتیاز منفی می‌دهد. اگر $p=P(T)$، کران نظری $p(1-p)$ در $S=T$ حاصل می‌شود؛ ممکن است هیچ توصیف قابل استفاده‌ای دقیقاً این گروه را بیان نکند. خود WRAcc ضدیکنوا نیست؛ کم بودن امتیاز یک مقدم، بدون کران معتبر، مجوز هرس همهٔ تخصصی‌سازی‌های آن نیست.

\begin{example}{کشف زیرگروه در مثال فیل}
شمارش‌ها $TP=1894,FP=0,FN=32,TN=54874$ و جمعیت $N=56800$ است. خلوص یک است و
\begin{align*}
P(T)&=1926/56800\approx0.03391,\\
P(S)&=1894/56800\approx0.03335,\\
\operatorname{WRAcc}&\approx0.03221.
\end{align*}
کران نظری $P(T)(1-P(T))\approx0.03276$ است. ابتدا از شمارش‌های دقیق استفاده کنید و در پایان گرد کنید.
\end{example}

\section{خودآزمایی و تمرین}
پیش از دیدن پاسخ، شمارش‌ها و استدلال خود را بنویسید.
\begin{selftests}
\begin{selftest}{شمارش با احتمال شرطی چه فرقی دارد؟}
در سبدهای مرجع، قاعدهٔ نان$\to$شیر را بررسی کنید. پشتیبانی قاعده، اطمینان و Lift را به دست آورید و بگویید آیا دیدن نان احتمال شیر را افزایش داده است؟
\answer{شمارش نان ۴، شیر ۴ و اجتماع ۳ است:
\[s=\frac35=0.6,\qquad c=\frac34=0.75,\qquad\operatorname{Lift}=\frac{0.75}{0.8}=0.9375.\]
پس با وجود اطمینان ۷۵ درصد، احتمال شیر از ۸۰ درصد به ۷۵ درصد کاهش یافته است. اطمینان را باید با نرخ پایه مقایسه کرد.}
\end{selftest}

\begin{selftest}{آیا همهٔ زوج‌های پرتکرار، سه‌تایی پرتکرار می‌سازند؟}
در دادهٔ مرجع و با حداقل شمارش ۳، مجموعه‌های پرتکرار هر سطح و نامزد سه‌تایی را پیدا کنید. آیا پرتکرار بودن تمام زوج‌ها کافی است؟
\answer{\begin{enumerate}
\item تک‌اقلام پرتکرار $F_1=\{B,M,D,R\}$ هستند؛ $C,E$ حذف می‌شوند.
\item شمارش شش زوج: $BM=3,BD=3,BR=2,MD=3,MR=2,DR=3$.
\item بنابراین $F_2=\{BM,BD,MD,DR\}$ است.
\item تنها نامزد سه‌تایی $BMD$ است، چون همهٔ زوج‌هایش پرتکرارند. شمارش واقعی آن ۲ است و حذف می‌شود.
\item $F_3$ تهی و فرایند تمام می‌شود. در این اجرا $6+6+1=13$ نامزد شمرده شده، درحالی‌که تعداد مجموعه‌های ناتهی تا اندازهٔ سه ۴۱ و تمام اندازه‌ها ۶۳ است.
\end{enumerate}


 $BMD$ فقط در تراکنش۴ و۵ ظاهر می‌شود، پس شمارش۲ دارد. اصل Apriori می‌گوید پرتکراری مجموعه، پرتکراری زیرمجموعه‌ها را لازم می‌کند؛ شرط کافی نمی‌دهد. هنگام تولید نامزد هنوز باید شمارش واقعی انجام شود.}
\end{selftest}

\begin{selftest}{آستانهٔ نسبی و گسستگی شمارش}
در۵ تراکنش اگر حداقل حمایت۵۰٪ باشد، چه حداقل شمارشی لازم است؟ آیا مجموعهٔ دارای شمارش۲ پذیرفته می‌شود؟
\answer{$\lceil5\times0.5\rceil=3$. شمارش۲ حمایت۴۰٪ دارد و رد می‌شود. آستانه با گردکردن معمولی به۲ تبدیل نمی‌شود.}
\end{selftest}

\begin{selftest}{تعداد قواعد}
از یک مجموعهٔ پرتکرار چهارعضوی چند قاعده می‌توان ساخت؟ چرا عدد۱۶ صحیح نیست؟
\answer{$2^4-2=14$؛ انتخاب مقدم تهی و مقدم کل مجموعه، یکی از دو سوی قاعده را تهی می‌کنند و کنار گذاشته می‌شوند. قواعد متفاوت می‌توانند پشتیبانی یکسان و اطمینان متفاوت داشته باشند.}
\end{selftest}

\begin{selftest}{هرس اعتماد با شرط صحیح}
از مجموعهٔ ثابت $L=\{A,B,C,D\}$، قاعدهٔ $ABC\to D$ اطمینان $0.6$ دارد و $\mathrm{minconf}=0.7$ است. دربارهٔ $AB\to CD$ و $A\to BCD$ چه می‌دانیم؟ آیا همین منطق برای $AB\to E$ هم قابل استفاده است؟
\answer{دو قاعدهٔ اول اطمینانی حداکثر $0.6$ دارند و قابل هرس‌اند؛ اجتماع اقلام هر سه قاعده $L$ ثابت است و کوچک‌کردن مقدم مخرج را کاهش نمی‌دهد. دربارهٔ $AB\to E$ چنین نتیجه‌ای نداریم، چون اجتماع عوض شده است.}
\end{selftest}

\begin{selftest}{معکوس قاعده}
در دادهٔ مرجع $c(D\to R)=0.75$ ولی $c(R\to D)=1$ است. آیا تناقض است؟ بالابری دو جهت را محاسبه کنید.
\answer{نه؛ جمعیت شرطی جهت اول ۴ و جهت دوم ۳ تراکنش است. بالابری در دو جهت برابر است:
\[\frac{0.75}{0.6}=1.25,\qquad\frac{1}{0.8}=1.25.\]
نرخ پایه تفاوت جهت را خنثی کرده است.}
\end{selftest}

\begin{selftest}{برداشتن یک قلم کم‌تکرار}
اگر $E$ فقط در یک تراکنش حضور دارد و حداقل شمارش۳ است، آیا لازم است همهٔ مجموعه‌های شامل $E$ را تک‌تک بشماریم؟
\answer{خیر؛ هر مجموعهٔ شامل $E$ شمارش حداکثر۱ دارد. با شش قلم، $2^5=32$ مجموعهٔ شامل $E$ وجود دارد و می‌توان آن‌ها را یکجا با اصل ضدیکنوایی حذف کرد (در شمارش اجرای خاص، برخی ممکن است از قبل به دلیل اقلام دیگر هرس شده باشند؛ تعداد صرفه‌جویی را دو بار نشماریم).}
\end{selftest}

\begin{selftest}{خلوص در برابر پوشش}
در جامعهٔ۱۰۰ نفری با۲۰ عضو هدف، گروه $S_1$ شامل۵ نفر است و هر۵ نفر هدف‌اند. گروه $S_2$ شامل۲۰ نفر است و۱۲ نفر هدف‌اند. کدام اطمینان و کدام WRAcc بیشتری دارد؟
\answer{گروه اول پوشش $0.05$ و خلوص یک دارد؛ بنابراین $\operatorname{WRAcc}=0.04$. گروه دوم پوشش $0.2$ و خلوص $0.6$ دارد و امتیازش $0.08$ است. اولی خالص‌تر و دومی از دید این معیار جذاب‌تر است؛ زیرا افزایش نرخ هدف را برای جمعیت بیشتری ایجاد می‌کند.}
\end{selftest}

\begin{selftest}{درخت FP و پشتیبانی}
در FP-tree قلم $A$ در دو گره با شمارنده‌های۳ و۲ ظاهر شده است. آیا می‌توان گفت پشتیبانی شمارشی $A$ برابر۳ است؟ برای ساخت پایگاه شرطی پسوند $A$ از وزن کدام گره استفاده می‌کنیم؟
\answer{شمارش کل $A$ جمع۵ است. هر مسیر پیشوندی باید با شمارندهٔ گره $A$ همان مسیر وزن بگیرد، نه شمارندهٔ والد یا ریشه. در غیر این صورت تراکنش‌های فاقد پسوند به اشتباه شمرده می‌شوند.}
\end{selftest}

\begin{selftest}{مجموعهٔ بسته و بیشینه}
فرض کنید شمارش‌های $\sigma(A)=5$، $\sigma(AB)=5$ و $\sigma(ABC)=3$ است و حداقل شمارش۳. آیا $A$ بسته است؟ آیا $AB$ بیشینه است؟
\answer{$A$ بسته نیست، چون $AB$ ابرمجموعه‌ای با همان شمارش دارد. $AB$ بیشینه نیست چون $ABC$ نیز پرتکرار است؛ ممکن است $AB$ بسته باشد، مشروط به اینکه هیچ ابرمجموعهٔ دیگر با شمارش۵ نداشته باشد. بسته بودن و بیشینه بودن دو سؤال متفاوت‌اند.}
\end{selftest}

\begin{selftest}{صفرهای مشترک}
جدول دو قلم $[a,b;c,d]=[70,10;10,10]$ است. ۱۰۰ تراکنش فاقد هر دو قلم اضافه کنیم. چه چیزی دربارهٔ Jaccard ، cosine و Lift تغییر می‌کند؟
\answer{مقادیر $J=70/90=7/9$ و $\operatorname{cos}=70/80=0.875$ ثابت‌اند. جمعیت $N$ از ۱۰۰ به ۲۰۰ می‌رسد و بالابری دو برابر می‌شود:
\begin{align*}
\operatorname{Lift}_{\mathrm{before}}&=\frac{100\times70}{80\times80}=1.09375,\\
\operatorname{Lift}_{\mathrm{after}}&=\frac{200\times70}{80\times80}=2.1875.
\end{align*}
پس انتخاب معیار با تعریف جامعه و اهمیت غیاب مشترک پیوند دارد.}
\end{selftest}

\begin{selftest}{چرا تغییر مرز گسسته‌سازی مهم است؟}
گروه سن۱۶--۲۴ در نمونهٔ۲۵۰ نفری،۲۲ پاسخ مثبت از۲۷ نفر دارد. اگر آن را به دو بازه تقسیم کنیم که هرکدام۱۱ پاسخ مثبت دارند، در حداقل حمایت۵٪ چه رخ می‌دهد؟
\answer{گروه اصلی حمایت۸٫۸٪ دارد؛ هر زیرقاعده حمایت۴٫۴٪ دارد و رد می‌شود، حتی اگر اطمینان هرکدام بالا باشد. گسسته‌سازی صرفاً آرایش نمایش نیست؛ فضای قواعد قابل کشف را تغییر می‌دهد.}
\end{selftest}
\end{selftests}
